\documentclass[a4paper, final]{article}

\usepackage[russian]{babel}
\usepackage{fontspec}
\setmainfont{Times New Roman}
\newfontfamily\cyrillicfonttt{Courier New}
\usepackage[14pt]{extsizes}
\usepackage[left=20mm, top=20mm, right=20mm, bottom=20mm, footskip=10mm]{geometry}
\usepackage{ragged2e}
\usepackage{indentfirst}
\usepackage{graphicx}
\usepackage{caption}
\usepackage{xcolor}
\usepackage{listings}
\usepackage{booktabs}
\usepackage{hyperref}
\usepackage{amsmath}
\usepackage{amssymb}
\usepackage{array}
\usepackage{multirow}

\definecolor{link_color}{HTML}{0000EE}
\hypersetup{colorlinks=true,linkcolor=link_color,urlcolor=link_color,citecolor=link_color}

\lstset{
    language=Python,
    basicstyle=\cyrillicfonttt\small,
    numbers=left,
    numberstyle=\tiny\color{gray},
    keywordstyle=\color{blue}\bfseries,
    commentstyle=\color{gray}\itshape,
    stringstyle=\color{red},
    frame=single,
    captionpos=b,
    breaklines=true,
    breakatwhitespace=false,
    showstringspaces=false,
    keepspaces=true,
    columns=flexible,
    fontadjust=true,
    basewidth={0.5em,0.45em},
    tabsize=4,
}

\begin{document}

{
\thispagestyle{empty}
\begin{center}
    \hfill\break\hfill\break
    МИНИСТЕРСТВО НАУКИ И ВЫСШЕГО ОБРАЗОВАНИЯ РОССИЙСКОЙ ФЕДЕРАЦИИ\\[5pt]
    Федеральное государственное автономное образовательное учреждение высшего образования
    \guillemotleft Санкт-Петербургский политехнический университет Петра Великого\guillemotright\\[10pt]
    Институт компьютерных наук и кибербезопасности\\[5pt]
    Направление: 02.03.01 Математика и компьютерные науки\\[30pt]
    {\large Отчет по дисциплине: \guillemotleft Вычислительная математика\guillemotright}\\[30pt]
    \LARGE\bf \guillemotleft Лабораторная работа 5\guillemotright.
\end{center}
\vspace{110pt}
\begin{tabular*}{460pt}{@{\extracolsep{\fill}} l r}
    Студент,\tabularnewline группы 5130201/40003 \hspace{50pt} & \line(1,0){100}\quad Адиатуллин Т.\,Р.\\
    & \\
    Руководитель,\tabularnewline Преподаватель & \line(1,0){121.5}\quad Пак В.\,Г.\\
    & \\ & \\ & \\
    & \guillemotleft\line(1,0){30}\guillemotright\line(1,0){100}\,2026\,г.
\end{tabular*}\\
\vfill
\begin{center}Санкт-Петербург, 2026\end{center}
\newpage
}

\tableofcontents
\newpage

\section{Постановка задачи}

\textbf{Тема:} численное решение линейных систем.

\textbf{Цель:} изучить прямые и итерационные методы решения систем линейных уравнений.

\textbf{Задание:}
\begin{enumerate}
    \item Исследовать линейную систему методом Гаусса. Найти решение.
    \item Решить систему итерационными методами: Якоби и Зейделя. Вычислить пять итераций, оценить погрешность.
\end{enumerate}

Дана система $Ax = b$ (вариант~1):
\[
A = \begin{pmatrix}
1{,}7   & 0{,}03  & 0{,}04  & 0{,}05 \\
0       & 0{,}8   & 0{,}01  & 0{,}02 \\
-0{,}03 & -0{,}02 & -0{,}1  & 0      \\
-0{,}05 & -0{,}04 & -0{,}03 & -1
\end{pmatrix}, \qquad
b = \begin{pmatrix}
0{,}681 \\ 0{,}4803 \\ -0{,}0802 \\ -1{,}0007
\end{pmatrix}.
\]

\newpage

\section{Задание 1. Метод Гаусса}

\subsection{Теоретические сведения}

Метод Гаусса~--- прямой метод решения системы $Ax = b$, состоящий из двух этапов.

\textbf{Прямой ход.} Расширенная матрица $[A\,|\,b]$ приводится к верхнетреугольному
виду. На шаге $k$ для каждой строки $i > k$ из строки $i$ вычитается строка $k$,
умноженная на множитель
\[
m_{ik} = \frac{a_{ik}^{(k)}}{a_{kk}^{(k)}}.
\]

\textbf{Обратный ход.} Из полученной треугольной системы находятся неизвестные:
\[
x_n = \frac{b_n^{(n)}}{a_{nn}^{(n)}}, \qquad
x_i = \frac{1}{a_{ii}^{(i)}}\!\left(b_i^{(i)} - \sum_{j=i+1}^{n} a_{ij}^{(i)}\,x_j\right),
\quad i = n-1, \ldots, 1.
\]

\subsection{Прямой ход}

Запишем расширенную матрицу $[A\,|\,b]$:
\[
\left(\begin{array}{cccc|c}
1{,}7   & 0{,}03  & 0{,}04  & 0{,}05 & 0{,}681  \\
0       & 0{,}8   & 0{,}01  & 0{,}02 & 0{,}4803 \\
-0{,}03 & -0{,}02 & -0{,}1  & 0      & -0{,}0802 \\
-0{,}05 & -0{,}04 & -0{,}03 & -1     & -1{,}0007
\end{array}\right).
\]

\textbf{Шаг 1.} Ведущий элемент $a_{11} = 1{,}7$. Множители:
\[
m_{31} = \frac{-0{,}03}{1{,}7} = -0{,}01765, \qquad
m_{41} = \frac{-0{,}05}{1{,}7} = -0{,}02941.
\]
($m_{21} = 0$, так как $a_{21} = 0$.) Выполняем $R_3 \leftarrow R_3 - m_{31} R_1$,
$R_4 \leftarrow R_4 - m_{41} R_1$:
\[
\left(\begin{array}{cccc|c}
1{,}7 & 0{,}03       & 0{,}04       & 0{,}05       & 0{,}681    \\
0     & 0{,}8        & 0{,}01       & 0{,}02       & 0{,}4803   \\
0     & -0{,}019471  & -0{,}099294  & 0{,}000882   & -0{,}068182 \\
0     & -0{,}039118  & -0{,}028824  & -0{,}998529  & -0{,}980671
\end{array}\right).
\]

\textbf{Шаг 2.} Ведущий элемент $a_{22}^{(2)} = 0{,}8$. Множители:
\[
m_{32} = \frac{-0{,}019471}{0{,}8} = -0{,}024338, \qquad
m_{42} = \frac{-0{,}039118}{0{,}8} = -0{,}048897.
\]
Выполняем $R_3 \leftarrow R_3 - m_{32} R_2$, $R_4 \leftarrow R_4 - m_{42} R_2$:
\[
\left(\begin{array}{cccc|c}
1{,}7 & 0{,}03 & 0{,}04       & 0{,}05       & 0{,}681    \\
0     & 0{,}8  & 0{,}01       & 0{,}02       & 0{,}4803   \\
0     & 0      & -0{,}099051  & 0{,}001369   & -0{,}056493 \\
0     & 0      & -0{,}028335  & -0{,}997551  & -0{,}957185
\end{array}\right).
\]

\textbf{Шаг 3.} Ведущий элемент $a_{33}^{(3)} = -0{,}099051$. Множитель:
\[
m_{43} = \frac{-0{,}028335}{-0{,}099051} = 0{,}286061.
\]
Выполняем $R_4 \leftarrow R_4 - m_{43} R_3$:
\[
\left(\begin{array}{cccc|c}
1{,}7 & 0{,}03 & 0{,}04       & 0{,}05       & 0{,}681    \\
0     & 0{,}8  & 0{,}01       & 0{,}02       & 0{,}4803   \\
0     & 0      & -0{,}099051  & 0{,}001369   & -0{,}056493 \\
0     & 0      & 0            & -0{,}997943  & -0{,}941025
\end{array}\right).
\]
Все диагональные элементы ненулевые~--- система имеет единственное решение.

\subsection{Обратный ход}

Из уравнения 4:
\[
x_4 = \frac{-0{,}941025}{-0{,}997943} = 0{,}94296.
\]

Из уравнения 3:
\[
x_3 = \frac{-0{,}056493 - 0{,}001369 \cdot 0{,}94296}{-0{,}099051} = \]
     \[ = \frac{-0{,}056493 - 0{,}001291}{-0{,}099051}
     = \frac{-0{,}057784}{-0{,}099051} = 0{,}58338.
\]

Из уравнения 2:
\[
x_2 = \frac{0{,}4803 - 0{,}01 \cdot 0{,}58338 - 0{,}02 \cdot 0{,}94296}{0{,}8} = \]
     \[ = \frac{0{,}4803 - 0{,}005834 - 0{,}018859}{0{,}8}
     = \frac{0{,}455607}{0{,}8} = 0{,}56951.
\]

Из уравнения 1:
\[
x_1 = \frac{0{,}681 - 0{,}03 \cdot 0{,}56951 - 0{,}04 \cdot 0{,}58338 - 0{,}05 \cdot 0{,}94296}{1{,}7}
\]
\[
= \frac{0{,}681 - 0{,}017085 - 0{,}023335 - 0{,}047148}{1{,}7}
= \frac{0{,}593432}{1{,}7} = 0{,}34908.
\]

\textbf{Решение:}
\[
\boxed{x^* = (0{,}34908;\quad 0{,}56951;\quad 0{,}58338;\quad 0{,}94296).}
\]

\textbf{Проверка} $Ax^* = b$:
\[1{,}7 \cdot 0{,}34908 + 0{,}03 \cdot 0{,}56951 + 0{,}04 \cdot 0{,}58338 + 0{,}05 \cdot 0{,}94296 = \] \[= 0{,}5934 + 0{,}01709 + 0{,}02334 + 0{,}04715 = 0{,}681\ \checkmark\]
\[0{,}8 \cdot 0{,}56951 + 0{,}01 \cdot 0{,}58338 + 0{,}02 \cdot 0{,}94296
 = \] \[ = 0{,}45561 + 0{,}00583 + 0{,}01886 = 0{,}4803\ \checkmark\]
\[-0{,}03 \cdot 0{,}34908 - 0{,}02 \cdot 0{,}56951 - 0{,}1 \cdot 0{,}58338
= \] \[= -0{,}01047 - 0{,}01139 - 0{,}05834 = -0{,}0802\ \checkmark\]
\[-0{,}05 \cdot 0{,}34908 - 0{,}04 \cdot 0{,}56951 - 0{,}03 \cdot 0{,}58338 - 1 \cdot 0{,}94296
= \] \[= -0{,}01745 - 0{,}02278 - 0{,}01750 - 0{,}94296 = -1{,}0007\ \checkmark\]

\newpage

\section{Задание 2. Итерационные методы}

\subsection{Приведение системы к виду $x = Bx + c$}

Разрешим каждое уравнение системы $Ax = b$ относительно
соответствующего неизвестного:
\[
\left\{
\begin{aligned}
x_1 &= -\frac{a_{12}}{a_{11}}x_2 - \frac{a_{13}}{a_{11}}x_3 - \frac{a_{14}}{a_{11}}x_4
        + \frac{b_1}{a_{11}},\\
x_2 &= -\frac{a_{23}}{a_{22}}x_3 - \frac{a_{24}}{a_{22}}x_4 + \frac{b_2}{a_{22}},\\
x_3 &= -\frac{a_{31}}{a_{33}}x_1 - \frac{a_{32}}{a_{33}}x_2 + \frac{b_3}{a_{33}},\\
x_4 &= -\frac{a_{41}}{a_{44}}x_1 - \frac{a_{42}}{a_{44}}x_2
        - \frac{a_{43}}{a_{44}}x_3 + \frac{b_4}{a_{44}}.
\end{aligned}
\right.
\]

Коэффициенты $b_{ij} = -a_{ij}/a_{ii}$ (при $i \ne j$) и $c_i = b_i/a_{ii}$:
\begin{align*}
b_{12} &= -\tfrac{0{,}03}{1{,}7}    = -0{,}01765, \\
b_{13} &= -\tfrac{0{,}04}{1{,}7}    = -0{,}02353, \\
b_{14} &= -\tfrac{0{,}05}{1{,}7}    = -0{,}02941.
\end{align*}

\begin{align*}
b_{23} &= -\tfrac{0{,}01}{0{,}8}    = -0{,}01250, \\
b_{24} &= -\tfrac{0{,}02}{0{,}8}    = -0{,}02500.
\end{align*}

\begin{align*}
b_{31} &= -\tfrac{-0{,}03}{-0{,}1}  = -0{,}30000, \\
b_{32} &= -\tfrac{-0{,}02}{-0{,}1}  = -0{,}20000.
\end{align*}

\begin{align*}
b_{41} &= -\tfrac{-0{,}05}{-1}      = -0{,}05000, \\
b_{42} &= -\tfrac{-0{,}04}{-1}      = -0{,}04000, \\
b_{43} &= -\tfrac{-0{,}03}{-1}      = -0{,}03000.
\end{align*}

\begin{align*}
	c_1    &= \tfrac{0{,}681}{1{,}7}     =  0{,}40059, \\
c_2    &= \tfrac{0{,}4803}{0{,}8}    =  0{,}60038, \\
c_3    &= \tfrac{-0{,}0802}{-0{,}1} =  0{,}80200, \\
c_4    &= \tfrac{-1{,}0007}{-1}     =  1{,}00070.
\end{align*}

Итерационная матрица $B$ и вектор $c$:
\[
B = \begin{pmatrix}
0          & -0{,}01765 & -0{,}02353 & -0{,}02941 \\
0          & 0          & -0{,}01250 & -0{,}02500 \\
-0{,}30000 & -0{,}20000 & 0          & 0          \\
-0{,}05000 & -0{,}04000 & -0{,}03000 & 0
\end{pmatrix},\qquad
c = \begin{pmatrix}0{,}40059\\0{,}60038\\0{,}80200\\1{,}00070\end{pmatrix}.
\]

\subsection{Условие сходимости}

Норма матрицы $B$ (максимальная сумма по строке~--- норма $\|\cdot\|_\infty$):
\begin{align*}
\text{строка 1:}&\quad 0{,}01765+0{,}02353+0{,}02941 = 0{,}07059,\\
\text{строка 2:}&\quad 0{,}01250+0{,}02500           = 0{,}03750,\\
\text{строка 3:}&\quad 0{,}30000+0{,}20000           = 0{,}50000,\\
\text{строка 4:}&\quad 0{,}05000+0{,}04000+0{,}03000 = 0{,}12000.
\end{align*}
\[
\|B\|_\infty = \max\{0{,}07059;\;0{,}03750;\;0{,}50000;\;0{,}12000\} = 0{,}5 < 1.
\]
По теореме~1 оба метода сходятся при любом начальном приближении.

\newpage

\subsection{Метод Якоби}

Расчётная формула метода Якоби~--- на каждой итерации используются
компоненты предыдущего приближения $x^{(k-1)}$:
\[
\left\{
\begin{aligned}
x_1^{(k)} &= b_{12}\,x_2^{(k-1)} + b_{13}\,x_3^{(k-1)} + b_{14}\,x_4^{(k-1)} + c_1,\\
x_2^{(k)} &= b_{23}\,x_3^{(k-1)} + b_{24}\,x_4^{(k-1)} + c_2,\\
x_3^{(k)} &= b_{31}\,x_1^{(k-1)} + b_{32}\,x_2^{(k-1)} + c_3,\\
x_4^{(k)} &= b_{41}\,x_1^{(k-1)} + b_{42}\,x_2^{(k-1)} + b_{43}\,x_3^{(k-1)} + c_4,
\end{aligned}
\right.\qquad k=1,2,\ldots
\]
Начальное приближение: $x^{(0)} = (0;\;0;\;0;\;0)^T$.

\medskip\textbf{Итерация $k=1$:}
\begin{align*}
x_1^{(1)} &= (-0{,}01765)\cdot0 + (-0{,}02353)\cdot0 + (-0{,}02941)\cdot0 + 0{,}40059 = \mathbf{0{,}40059},\\
x_2^{(1)} &= (-0{,}01250)\cdot0 + (-0{,}02500)\cdot0 + 0{,}60038 = \mathbf{0{,}60038},\\
x_3^{(1)} &= (-0{,}30000)\cdot0 + (-0{,}20000)\cdot0 + 0{,}80200 = \mathbf{0{,}80200},\\
x_4^{(1)} &= (-0{,}05000)\cdot0 + (-0{,}04000)\cdot0 + (-0{,}03000)\cdot0 + 1{,}00070 = \mathbf{1{,}00070}.
\end{align*}
$\|x^{(1)}-x^{(0)}\|_\infty = \max\{0{,}40059;\;0{,}60038;\;0{,}80200;\;1{,}00070\} = 1{,}00070.$

\medskip\textbf{Итерация $k=2$:}
\begin{align*}
x_1^{(2)} &= (-0{,}01765)\cdot0{,}60038 + (-0{,}02353)\cdot0{,}80200 + (-0{,}02941)\cdot1{,}00070 + 0{,}40059\\
           &= -0{,}01060 - 0{,}01887 - 0{,}02943 + 0{,}40059 = \mathbf{0{,}34169},\\
x_2^{(2)} &= (-0{,}01250)\cdot0{,}80200 + (-0{,}02500)\cdot1{,}00070 + 0{,}60038\\
           &= -0{,}01003 - 0{,}02502 + 0{,}60038 = \mathbf{0{,}56533},\\
x_3^{(2)} &= (-0{,}30000)\cdot0{,}40059 + (-0{,}20000)\cdot0{,}60038 + 0{,}80200\\
           &= -0{,}12018 - 0{,}12008 + 0{,}80200 = \mathbf{0{,}56175},\\
x_4^{(2)} &= (-0{,}05000)\cdot0{,}40059 + (-0{,}04000)\cdot0{,}60038 + (-0{,}03000)\cdot0{,}80200 + 1{,}00070\\
           &= -0{,}02003 - 0{,}02402 - 0{,}02406 + 1{,}00070 = \mathbf{0{,}93260}.
\end{align*}
$\|x^{(2)}-x^{(1)}\|_\infty = \max\{0{,}05890;\;0{,}03505;\;0{,}24025;\;0{,}06810\} = 0{,}24025.$

\medskip\textbf{Итерация $k=3$:}
\begin{align*}
x_1^{(3)} &= (-0{,}01765)\cdot0{,}56533 + (-0{,}02353)\cdot0{,}56175 + (-0{,}02941)\cdot0{,}93260 + 0{,}40059\\
           &= -0{,}00998 - 0{,}01322 - 0{,}02743 + 0{,}40059 = \mathbf{0{,}34997},\\
x_2^{(3)} &= (-0{,}01250)\cdot0{,}56175 + (-0{,}02500)\cdot0{,}93260 + 0{,}60038\\
           &= -0{,}00702 - 0{,}02332 + 0{,}60038 = \mathbf{0{,}57004},\\
x_3^{(3)} &= (-0{,}30000)\cdot0{,}34169 + (-0{,}20000)\cdot0{,}56533 + 0{,}80200\\
           &= -0{,}10251 - 0{,}11307 + 0{,}80200 = \mathbf{0{,}58643},\\
x_4^{(3)} &= (-0{,}05000)\cdot0{,}34169 + (-0{,}04000)\cdot0{,}56533 + (-0{,}03000)\cdot0{,}56175 + 1{,}00070\\
           &= -0{,}01708 - 0{,}02261 - 0{,}01685 + 1{,}00070 = \mathbf{0{,}94415}.
\end{align*}
$\|x^{(3)}-x^{(2)}\|_\infty = \max\{0{,}00828;\;0{,}00471;\;0{,}02468;\;0{,}01155\} = 0{,}02468.$

\medskip\textbf{Итерация $k=4$:}
\begin{align*}
x_1^{(4)} &= (-0{,}01765)\cdot0{,}57004 + (-0{,}02353)\cdot0{,}58643 + (-0{,}02941)\cdot0{,}94415 + 0{,}40059\\
           &= -0{,}01006 - 0{,}01380 - 0{,}02777 + 0{,}40059 = \mathbf{0{,}34896},\\
x_2^{(4)} &= (-0{,}01250)\cdot0{,}58643 + (-0{,}02500)\cdot0{,}94415 + 0{,}60038\\
           &= -0{,}00733 - 0{,}02360 + 0{,}60038 = \mathbf{0{,}56944},\\
x_3^{(4)} &= (-0{,}30000)\cdot0{,}34997 + (-0{,}20000)\cdot0{,}57004 + 0{,}80200\\
           &= -0{,}10499 - 0{,}11401 + 0{,}80200 = \mathbf{0{,}58300},\\
x_4^{(4)} &= (-0{,}05000)\cdot0{,}34997 + (-0{,}04000)\cdot0{,}57004 + (-0{,}03000)\cdot0{,}58643 + 1{,}00070\\
           &= -0{,}01750 - 0{,}02280 - 0{,}01759 + 1{,}00070 = \mathbf{0{,}94281}.
\end{align*}
$\|x^{(4)}-x^{(3)}\|_\infty = \max\{0{,}00101;\;0{,}00060;\;0{,}00343;\;0{,}00134\} = 0{,}00343.$

\medskip\textbf{Итерация $k=5$:}
\begin{align*}
x_1^{(5)} &= (-0{,}01765)\cdot0{,}56944 + (-0{,}02353)\cdot0{,}58300 + (-0{,}02941)\cdot0{,}94281 + 0{,}40059\\
           &= -0{,}01005 - 0{,}01372 - 0{,}02773 + 0{,}40059 = \mathbf{0{,}34909},\\
x_2^{(5)} &= (-0{,}01250)\cdot0{,}58300 + (-0{,}02500)\cdot0{,}94281 + 0{,}60038\\
           &= -0{,}00729 - 0{,}02357 + 0{,}60038 = \mathbf{0{,}56952},\\
x_3^{(5)} &= (-0{,}30000)\cdot0{,}34896 + (-0{,}20000)\cdot0{,}56944 + 0{,}80200\\
           &= -0{,}10469 - 0{,}11389 + 0{,}80200 = \mathbf{0{,}58342},\\
x_4^{(5)} &= (-0{,}05000)\cdot0{,}34896 + (-0{,}04000)\cdot0{,}56944 + (-0{,}03000)\cdot0{,}58300 + 1{,}00070\\
           &= -0{,}01745 - 0{,}02278 - 0{,}01749 + 1{,}00070 = \mathbf{0{,}94298}.
\end{align*}
$\|x^{(5)}-x^{(4)}\|_\infty = \max\{0{,}00013;\;0{,}00008;\;0{,}00042;\;0{,}00017\} = 0{,}00042.$

\begin{table}[h!]
\centering
\caption{Итерации метода Якоби}
\label{tab:jacobi}
\renewcommand{\arraystretch}{1.3}
\begin{tabular}{|c|c|c|c|c|c|c|}
\hline
 & $x^{(0)}$ & $x^{(1)}$ & $x^{(2)}$ & $x^{(3)}$ & $x^{(4)}$ & $x^{(5)}$ \\
\hline
$x_1$ & 0 & 0,40059 & 0,34169 & 0,34997 & 0,34896 & 0,34909 \\
$x_2$ & 0 & 0,60038 & 0,56533 & 0,57004 & 0,56944 & 0,56952 \\
$x_3$ & 0 & 0,80200 & 0,56175 & 0,58643 & 0,58300 & 0,58342 \\
$x_4$ & 0 & 1,00070 & 0,93260 & 0,94415 & 0,94281 & 0,94298 \\
\hline
$\|x^{(k)}-x^{(k-1)}\|_\infty$ & --- & 1,00070 & 0,24025 & 0,02468 & 0,00343 & 0,00042 \\
\hline
\end{tabular}
\end{table}

Апостериорная оценка погрешности (теорема~2):
\[
\Delta x^{(5)} \leq \frac{\|B\|}{1-\|B\|}\cdot\|x^{(5)}-x^{(4)}\|_\infty
= \frac{0{,}5}{0{,}5}\cdot 0{,}00042 = 0{,}00042.
\]
Фактическая погрешность: $\|x^{(5)}-x^*\|_\infty \approx 0{,}00005.$

\newpage

\subsection{Метод Зейделя}

В методе Зейделя при вычислении $x_i^{(k)}$ используются
уже найденные на шаге $k$ значения $x_1^{(k)},\ldots,x_{i-1}^{(k)}$.

Матрица $B$ разбивается: $B = B_1 + B_2$, где
\[
B_1 = \begin{pmatrix}
0          & 0          & 0          & 0 \\
0          & 0          & 0          & 0 \\
-0{,}30000 & -0{,}20000 & 0          & 0 \\
-0{,}05000 & -0{,}04000 & -0{,}03000 & 0
\end{pmatrix}, \] 
\[
B_2 = \begin{pmatrix}
0 & -0{,}01765 & -0{,}02353 & -0{,}02941 \\
0 & 0          & -0{,}01250 & -0{,}02500 \\
0 & 0          & 0          & 0          \\
0 & 0          & 0          & 0
\end{pmatrix}.
\]
\[
\|B_1\|_\infty = 0{,}500,\qquad \|B_2\|_\infty = 0{,}07059,\qquad
\|B_1\|_\infty+\|B_2\|_\infty = 0{,}57059 < 1.
\]
По теореме~3 метод сходится. Скорость сходимости определяется нормой $B_2$:
\[
q = \frac{\|B_2\|}{1-\|B_1\|} = \frac{0{,}07059}{1-0{,}500} = 0{,}14118.
\]

Расчётная формула метода Зейделя:
\[
\left\{
\begin{aligned}
x_1^{(k)} &= b_{12}\,x_2^{(k-1)} + b_{13}\,x_3^{(k-1)} + b_{14}\,x_4^{(k-1)} + c_1,\\
x_2^{(k)} &= b_{23}\,x_3^{(k-1)} + b_{24}\,x_4^{(k-1)} + c_2,\\
x_3^{(k)} &= b_{31}\,x_1^{(k)} + b_{32}\,x_2^{(k)} + c_3,\\
x_4^{(k)} &= b_{41}\,x_1^{(k)} + b_{42}\,x_2^{(k)} + b_{43}\,x_3^{(k)} + c_4.
\end{aligned}
\right.\qquad k=1,2,\ldots
\]
Начальное приближение: $x^{(0)} = (0;\;0;\;0;\;0)^T$.

\medskip\textbf{Итерация $k=1$:}
\begin{align*}
x_1^{(1)} &= (-0{,}01765)\cdot0 + (-0{,}02353)\cdot0 + (-0{,}02941)\cdot0 + 0{,}40059 = \mathbf{0{,}40059},\\
x_2^{(1)} &= (-0{,}01250)\cdot0 + (-0{,}02500)\cdot0 + 0{,}60038 = \mathbf{0{,}60038},\\
x_3^{(1)} &= (-0{,}30000)\cdot\mathbf{0{,}40059} + (-0{,}20000)\cdot\mathbf{0{,}60038} + 0{,}80200\\
           &= -0{,}12018 - 0{,}12008 + 0{,}80200 = \mathbf{0{,}56175},\\
x_4^{(1)} &= (-0{,}05000)\cdot\mathbf{0{,}40059} + (-0{,}04000)\cdot\mathbf{0{,}60038}
             + (-0{,}03000)\cdot\mathbf{0{,}56175} + 1{,}00070\\
           &= -0{,}02003 - 0{,}02402 - 0{,}01685 + 1{,}00070 = \mathbf{0{,}93980}.
\end{align*}
$\|x^{(1)}-x^{(0)}\|_\infty = 0{,}93980.$

\medskip\textbf{Итерация $k=2$:}
\begin{align*}
x_1^{(2)} &= (-0{,}01765)\cdot0{,}60038 + (-0{,}02353)\cdot0{,}56175 + (-0{,}02941)\cdot0{,}93980 + 0{,}40059\\
           &= -0{,}01060 - 0{,}01322 - 0{,}02764 + 0{,}40059 = \mathbf{0{,}34914},\\
x_2^{(2)} &= (-0{,}01250)\cdot0{,}56175 + (-0{,}02500)\cdot0{,}93980 + 0{,}60038\\
           &= -0{,}00702 - 0{,}02350 + 0{,}60038 = \mathbf{0{,}56986},\\
x_3^{(2)} &= (-0{,}30000)\cdot\mathbf{0{,}34914} + (-0{,}20000)\cdot\mathbf{0{,}56986} + 0{,}80200\\
           &= -0{,}10474 - 0{,}11397 + 0{,}80200 = \mathbf{0{,}58329},\\
x_4^{(2)} &= (-0{,}05000)\cdot\mathbf{0{,}34914} + (-0{,}04000)\cdot\mathbf{0{,}56986}
             + (-0{,}03000)\cdot\mathbf{0{,}58329} + 1{,}00070\\
           &= -0{,}01746 - 0{,}02279 - 0{,}01750 + 1{,}00070 = \mathbf{0{,}94295}.
\end{align*}
$\|x^{(2)}-x^{(1)}\|_\infty = \max\{0{,}05145;\;0{,}03052;\;0{,}02154;\;0{,}00315\} = 0{,}05145.$

\medskip\textbf{Итерация $k=3$:}
\begin{align*}
x_1^{(3)} &= (-0{,}01765)\cdot0{,}56986 + (-0{,}02353)\cdot0{,}58329 + (-0{,}02941)\cdot0{,}94295 + 0{,}40059\\
           &= -0{,}01006 - 0{,}01373 - 0{,}02773 + 0{,}40059 = \mathbf{0{,}34907},\\
x_2^{(3)} &= (-0{,}01250)\cdot0{,}58329 + (-0{,}02500)\cdot0{,}94295 + 0{,}60038\\
           &= -0{,}00729 - 0{,}02357 + 0{,}60038 = \mathbf{0{,}56951},\\
x_3^{(3)} &= (-0{,}30000)\cdot\mathbf{0{,}34907} + (-0{,}20000)\cdot\mathbf{0{,}56951} + 0{,}80200\\
           &= -0{,}10472 - 0{,}11390 + 0{,}80200 = \mathbf{0{,}58338},\\
x_4^{(3)} &= (-0{,}05000)\cdot\mathbf{0{,}34907} + (-0{,}04000)\cdot\mathbf{0{,}56951}
             + (-0{,}03000)\cdot\mathbf{0{,}58338} + 1{,}00070\\
           &= -0{,}01745 - 0{,}02278 - 0{,}01750 + 1{,}00070 = \mathbf{0{,}94297}.
\end{align*}
$\|x^{(3)}-x^{(2)}\|_\infty = \max\{0{,}00007;\;0{,}00035;\;0{,}00009;\;0{,}00002\} = 0{,}00035.$

\medskip\textbf{Итерации $k=4,5$} практически совпадают с точным решением
(изменения $< 10^{-5}$):
\[
x^{(4)} \approx x^{(5)} \approx (0{,}34908;\;0{,}56951;\;0{,}58338;\;0{,}94296)^T.
\]

\begin{table}[h!]
\centering
\caption{Итерации метода Зейделя}
\label{tab:seidel}
\renewcommand{\arraystretch}{1.3}
\begin{tabular}{|c|c|c|c|c|c|c|}
\hline
 & $x^{(0)}$ & $x^{(1)}$ & $x^{(2)}$ & $x^{(3)}$ & $x^{(4)}$ & $x^{(5)}$ \\
\hline
$x_1$ & 0 & 0,40059 & 0,34914 & 0,34907 & 0,34908 & 0,34908 \\
$x_2$ & 0 & 0,60038 & 0,56986 & 0,56951 & 0,56951 & 0,56951 \\
$x_3$ & 0 & 0,56175 & 0,58329 & 0,58338 & 0,58338 & 0,58338 \\
$x_4$ & 0 & 0,93980 & 0,94295 & 0,94297 & 0,94296 & 0,94296 \\
\hline
$\|x^{(k)}-x^{(k-1)}\|_\infty$ & --- & 0,93980 & 0,05145 & 0,00035 & $<10^{-4}$ & $<10^{-5}$ \\
\hline
\end{tabular}
\end{table}

Апостериорная оценка погрешности (теорема~4):
\[
\Delta x^{(3)} \leq \frac{\|B_2\|}{1-\|B\|}\cdot\|x^{(3)}-x^{(2)}\|_\infty
= \frac{0{,}07059}{0{,}5}\cdot 0{,}00035 = 0{,}14118\cdot 0{,}00035 = 0{,}00005.
\]
Фактическая погрешность: $\|x^{(3)}-x^*\|_\infty \approx 0{,}000004.$

\newpage

\section{Сравнение методов}

\begin{table}[h!]
\centering
\caption{Сравнение итерационных методов}
\label{tab:compare}
\renewcommand{\arraystretch}{1.4}
\begin{tabular}{|l|c|c|}
\hline
Показатель & Якоби & Зейделя \\
\hline
$\|B\|_\infty$ & $0{,}5$ & $0{,}5$ \\
\hline
$\|B_1\|_\infty + \|B_2\|_\infty$ & --- & $0{,}57059 < 1$ \\
\hline
$\|x^{(3)}-x^*\|_\infty$ & $0{,}003051$ & $0{,}000004$ \\
\hline
$\|x^{(5)}-x^*\|_\infty$ & $0{,}000048$ & $<10^{-6}$ \\
\hline
\end{tabular}
\end{table}

Метод Зейделя сходится значительно быстрее метода Якоби: уже после
3~итераций погрешность составляет $4\cdot10^{-6}$, тогда как у Якоби~---
$3\cdot10^{-3}$. Это объясняется тем, что при вычислении каждой компоненты
$x_i^{(k)}$ немедленно используются уточнённые значения предыдущих компонент.

\newpage

\section*{Приложение}
\addcontentsline{toc}{section}{Приложение}

\begin{lstlisting}[caption={Метод Гаусса и итерационные методы}, label=lst:code]
import numpy as np

A = np.array([
    [ 1.7,   0.03,  0.04,  0.05],
    [ 0.0,   0.8,   0.01,  0.02],
    [-0.03, -0.02, -0.1,   0.0 ],
    [-0.05, -0.04, -0.03, -1.0 ]
])
b = np.array([0.681, 0.4803, -0.0802, -1.0007])
n = 4

# ===== Metod Gaussa =====
Ab = np.hstack([A.astype(float), b.reshape(-1,1)])
for k in range(n):
    for i in range(k+1, n):
        m = Ab[i,k] / Ab[k,k]
        Ab[i] -= m * Ab[k]
x = np.zeros(n)
for i in range(n-1, -1, -1):
    x[i] = (Ab[i,n] - Ab[i,i+1:n] @ x[i+1:n]) / Ab[i,i]
print("Gauss:", np.round(x, 5))

# ===== Matritsa B i vektor c =====
B = np.zeros((n,n))
c = np.zeros(n)
for i in range(n):
    c[i] = b[i] / A[i,i]
    for j in range(n):
        if i != j:
            B[i,j] = -A[i,j] / A[i,i]
print("||B||_inf =", np.max(np.sum(np.abs(B), axis=1)))

# ===== Metod Yakobi =====
x_j = np.zeros(n)
for k in range(1, 6):
    x_j_new = B @ x_j  c
    diff = np.max(np.abs(x_j_new - x_j))
    print(f"Jacobi k={k}: {np.round(x_j_new,5)}, diff={diff:.5f}")
    x_j = x_j_new

# ===== Metod Zeydelya =====
x_s = np.zeros(n)
for k in range(1, 6):
    x_s_new = x_s.copy()
    for i in range(n):
        x_s_new[i] = c[i]
        for j in range(i):
            x_s_new[i] += B[i,j] * x_s_new[j]  # novye znacheniya
        for j in range(i+1, n):
            x_s_new[i] += B[i,j] * x_s[j]       # starye znacheniya
    diff = np.max(np.abs(x_s_new - x_s))
    print(f"Seidel k={k}: {np.round(x_s_new,5)}, diff={diff:.5f}")
    x_s = x_s_new
\end{lstlisting}

\end{document}